% !TEX root = vkr.tex
\section{Постановка задачи}

Целью работы является расширение возможностей библиотеки PySATL-CPD
для обработки последовательно поступающих данных посредством переноса
оконного компонента (скрабера) и подготовка воспроизводимой сравнительной
оценки выбранных алгоритмов обнаружения разладок.

Для достижения цели поставлены следующие задачи:
\begin{enumerate}
    \item Перенести скрабер из предыдущей версии PySATL-CPD
    в текущую архитектуру библиотеки:
    \begin{enumerate}
        \item адаптировать взаимодействие скрабера с поставщиками данных
        и совместимыми офлайн-алгоритмами;
        \item реализовать формирование скользящих окон
        и сопоставление локальных индексов окна с индексами исходного ряда;
        \item проверить корректность работы модульными тестами,
        включая обработку коротких последовательностей, перекрытие окон
        и завершение потока.
    \end{enumerate}

    \item Подобрать наборы данных для проверки алгоритмов
    при последовательной обработке наблюдений: описать источники
    и разметку разладок, подготовить загрузку и предобработку данных.
    При необходимости дополнить их синтетическими рядами
    с известными моментами изменений.

    \item Провести сравнительную оценку выбранных алгоритмов:
    \begin{enumerate}
        \item зафиксировать условия запуска и правила сопоставления
        обнаруженных и размеченных разладок;
        \item сравнить качество обнаружения, ложные срабатывания,
        задержку обнаружения и время обработки;
        \item сохранить параметры и результаты экспериментов
        для воспроизведения, сформулировать выводы о применимости методов
        и целесообразности расширения их набора.
    \end{enumerate}

    \item Подготовить пример использования перенесённого скрабера
    и документацию по его применению и воспроизведению экспериментов.
\end{enumerate}

Работа рассчитана на один семестр. В первую очередь планируется
перенос скрабера, затем подготовка данных и проведение экспериментов.
Документация дополняется по мере выполнения задач.

